% Purpose: Slides for the talk "Exploring Polymorphic Gradual Types with
% Phil and Peter".  This initial deck presents the reduction semantics from
% strong/Reduction.agda in a named-variable presentation, using the visual and
% notational style of old-slides.tex.

\documentclass[12pt]{beamer}
\usecolortheme{dove}
\usepackage[T1]{fontenc}
\usepackage{garamond}
\usefonttheme{serif}
\usepackage{graphicx}
\usepackage{amsmath}
\usepackage{amssymb}
\usepackage{semantic}
\usepackage{color}

\definecolor{darkgray}{gray}{0.6}
\definecolor{brightyellow}{rgb}{1,0.94,0}
\newcommand{\gray}[1]{{\color{darkgray} #1}}
\newcommand{\HI}[1]{\colorbox{brightyellow}{\color{black}\ensuremath{#1}}}
\newcommand{\kw}[1]{\mathtt{#1}}
\newcommand{\dotspace}{.\,}
\newcommand{\lam}[1]{\lambda #1 \dotspace}
\newcommand{\Lam}[1]{\Lambda #1 \dotspace}
\newcommand{\of}{{:}}
\newcommand{\by}{\mapsto}
\newcommand{\func}{\to}
\newcommand{\reduce}{\longrightarrow}
\newcommand{\boundary}[3]{\left[#1\right]^{#2}_{#3}}
\newcommand{\tyapp}[2]{#1\,#2}
\newcommand{\app}[2]{#1\,#2}
\newcommand{\seal}[1]{-#1}
\newcommand{\unseal}[1]{+#1}
\newcommand{\unsealat}[2]{\mathsf{reveal}(#1, #2)}
\newcommand{\unsealatc}[2]{\mathsf{reveal}(#1,#2)}
\newcommand{\idconv}[1]{\kw{id}\,#1}
\newcommand{\bind}[2]{#1{:=}#2}
\newcommand{\maskentry}[1]{\mathsf{mask}(#1)}
\newcommand{\morph}[2]{\left\langle #1\mathbin{;}#2\right\rangle}
\newcommand{\lockX}[1]{-#1}
\newcommand{\unlockX}[1]{+#1}
\newcommand{\dual}[1]{#1^{\mathrm{op}}}
\newcommand{\changes}[1]{\mathsf{changes}(#1)}
\newcommand{\unlocked}[1]{#1^{+}}
\newcommand{\scopemove}{\mathbin{\triangleleft}}
\newcommand{\mask}[2]{\mathsf{mask}_{#1}(#2)}
\newcommand{\unmask}[2]{\mathsf{unmask}_{#1}(#2)}
\newcommand{\pushbinds}[2]{\mathsf{pushBinds}(#1,#2)}
\newcommand{\hide}[1]{\mathsf{hide}(#1)}
\newcommand{\dualscope}[1]{\mathsf{dual}(#1)}
\newcommand{\rewind}[1]{\mathsf{rewind}(#1)}
\newcommand{\append}{\mathbin{+\!+}}
\newcommand{\idc}[1]{\mathsf{Id}(#1)}
\newcommand{\Value}[1]{\mathsf{Value}(#1)}
\newcommand{\entryname}[1]{\mathsf{name}(#1)}
\newcommand{\Nameable}[1]{\mathsf{Nameable}(#1)}
\newcommand{\Locked}[1]{\mathsf{Locked}(#1)}
\newcommand{\fce}[2]{\mathsf{convCtx}(#1,#2)}
\newcommand{\intc}[2]{#1(#2)}

\setbeamertemplate{navigation symbols}{}
\setbeamertemplate{footline}[frame number]
\setbeamercolor{footnote mark}{fg=white}

\garamond

\title[Exploring Polymorphic Gradual Types]{
  Exploring Polymorphic Gradual Types\\
  with Phil and Peter}
\author{Jeremy G. Siek}
\institute{Indiana University, Bloomington}
\date{}

\begin{document}

%===============================================================================
\frame{
\titlepage

\vspace{-30pt}
\begin{center}
\includegraphics[height=1.7in]{iu-tulips}
\end{center}
}
%===============================================================================
\frame{
  \frametitle{Terms, Conversions, and Context Morphisms}
\begingroup\small
\[
\begin{array}{lcl}
 A,B,C,D &::=& \iota \mid A \to B \mid X \mid \forall X.A \\
 L,M,N&::=& k \mid x \mid \lam{x\of A}N \mid \app{L}{M} \\
      && \mid \Lambda X. N \mid \tyapp{L}{A} \mid \HI{\boundary{M}{\Delta}{c}} \\
 \HI{c,d} &::=& \HI{\idconv{A} \mid \seal{X} \mid \unseal{X}
          \mid c \func d \mid \forall X.\,c} \\
 \HI{\delta} &::=& \HI{\emptyset \mid \lockX{X},\delta \mid \unlockX{X},\delta} \\
 \HI{\Theta} &::=& \HI{\emptyset \mid \bind{X}{A},\Theta} \\
 \HI{\Delta} &::=& \HI{\morph{\Theta}{\delta}}
\end{array}
\]
\endgroup
}
%===============================================================================
\frame{
  \frametitle{Contexts and type-variable lookup}


\[
\begin{array}{lcl}
 \psi &::=& X \mid \bind{X}{A} \mid \maskentry{\psi} \\
 \Psi &::=& \emptyset \mid \Psi,\psi
\end{array}
\]

\[
\entryname{X}=X
\qquad
\entryname{\bind{X}{A}}=X
\qquad
\entryname{\maskentry{\psi}}=\entryname{\psi}
\]

\begin{center}
\fbox{$\Psi\ni_{\!e}X,\psi$}
\end{center}

\vspace{-8pt}

{\small
\begin{columns}[T]
\begin{column}{0.48\textwidth}
\[
\begin{gathered}
\inference{}{\Psi,\psi\ni_{\!e}X,\psi} \\
\text{if }\entryname{\psi}=X
\end{gathered}
\]
\end{column}
\begin{column}{0.48\textwidth}
\[
\begin{gathered}
\inference{\Psi\ni_{\!e}X,\psi}{\Psi,\psi'\ni_{\!e}X,\psi} \\
\text{if }X\neq\entryname{\psi'}
\end{gathered}
\]
\end{column}
\end{columns}
}

\vspace{-2pt}

\hfill\fbox{$\Psi\ni X$}

\[
\inference{\Psi\ni_{\!e}X,\psi\qquad\Nameable{\psi}}
          {\Psi\ni X}
\qquad
\begin{array}{c}
\Nameable{X} \\
\Nameable{\bind{X}{A}}
\end{array}
\]
}
%===============================================================================
\frame{
  \frametitle{Locked type-variable lookup}

\[
\fbox{$\Psi\ni_{\mathrm{lk}}X$}
\]

\[
\inference{\Psi\ni_{\!e}X,\psi\qquad\Locked{\psi}}
          {\Psi\ni_{\mathrm{lk}}X}
\]

\[
\inference{\Nameable{\psi}}
          {\Locked{\maskentry{\psi}}}
\]
}
%===============================================================================
\frame{
  \frametitle{Operations on type contexts}

{\fontsize{8}{9}\selectfont
\[
\begin{aligned}
\mask{X}{\emptyset} &= \emptyset \\
\mask{X}{\Psi,\psi} &= \Psi,\maskentry{\psi}
&&\text{if }\entryname{\psi}=X \\
\mask{X}{\Psi,\psi} &= \mask{X}{\Psi},\psi
&&\text{if }X\neq\entryname{\psi}
\\[1.5ex]
\unmask{X}{\emptyset} &= \emptyset \\
\unmask{X}{\Psi,\maskentry{\psi}} &= \Psi,\psi
&&\text{if }\entryname{\psi}=X \\
\unmask{X}{\Psi,\psi} &= \Psi,\psi
&&\text{if }\entryname{\psi}=X\text{ and }\Nameable{\psi} \\
\unmask{X}{\Psi,\psi} &= \unmask{X}{\Psi},\psi
&&\text{if }X\neq\entryname{\psi}
\end{aligned}
\]
\vspace{-4pt}
\[
\begin{aligned}
\pushbinds{\emptyset}{\Psi}
  &= \Psi \\
\pushbinds{\bind{X}{A},\Theta}{\Psi}
  &= \pushbinds{\Theta}{\Psi},\bind{X}{A}
\end{aligned}
\]
}
}
%===============================================================================
\frame{
  \frametitle{Context morphisms acting on type contexts}

{\scriptsize
\begin{columns}[T]
\begin{column}{0.5\textwidth}
  \hfill\fbox{$\delta(\Psi) = \Psi'$}
  \[
\begin{aligned}
\emptyset(\Psi) &= \Psi \\
(\lockX{X},\delta)(\Psi) &= \mask{X}{\delta(\Psi)} \\
(\unlockX{X},\delta)(\Psi) &= \unmask{X}{\delta(\Psi)}
\end{aligned}
\]
\end{column}
\begin{column}{0.5\textwidth}
  \hfill\fbox{$\delta^{+}(\Psi) = \Psi'$}
\[
\begin{aligned}
\emptyset^{+}(\Psi) &= \Psi \\
(\lockX{X},\delta)^{+}(\Psi) &= \delta^{+}(\Psi) \\
(\unlockX{X},\delta)^{+}(\Psi) &= \unmask{X}{\delta^{+}(\Psi)}
\end{aligned}
\]
\end{column}
\end{columns}
\vspace{10pt}
 \hfill\fbox{$\Delta(\Psi) = \Psi'$}
\[
\begin{aligned}
\intc{\morph{\Theta}{\delta}}{\Psi}
  &= \pushbinds{\Theta}{\delta(\Psi)} \\
\fce{\morph{\Theta}{\delta}}{\Psi}
  &= \pushbinds{\Theta}{\delta^{+}(\Psi)}
\end{aligned}
\]
}
}
%===============================================================================
\frame{
  \frametitle{Dual, rewind, and scope move}

{\fontsize{7}{8}\selectfont
\begin{columns}[T]
\begin{column}{0.54\textwidth}
\[
\begin{aligned}
\hide{\emptyset} &= \emptyset \\
\hide{\bind{X}{A},\Theta} &= \lockX{X},\hide{\Theta} \\
\dualscope{\emptyset} &= \emptyset \\
\dualscope{\unlockX{X},\delta} &= \dualscope{\delta},\lockX{X} \\
\dualscope{\lockX{X},\delta} &= \dualscope{\delta},\unlockX{X} \\
\dual{\morph{\Theta}{\delta}}
  &= \morph{\emptyset}{\hide{\Theta}\append\dualscope{\delta}}
\end{aligned}
\]
\end{column}
\begin{column}{0.45\textwidth}
\[
\begin{aligned}
\rewind{\morph{\Theta}{\delta}}
  &= \morph{\Theta}{\dualscope{\delta}\append\delta} \\
\morph{\Theta_1}{\delta_1}\scopemove\morph{\Theta_2}{\delta_2}
  &= \morph{\Theta_1}{\delta_1\append\delta_2}
\end{aligned}
\]
\end{column}
\end{columns}
}
}
%===============================================================================
\frame{
  \frametitle{Conversion typing}

\hfill\fbox{$\Psi\vdash c:A\Rightarrow B$}

\begin{gather*}
\inference{}
          {\Psi\vdash\idconv{\iota}:\iota\Rightarrow\iota}
\quad
\inference{\Psi\ni X}
          {\Psi\vdash\idconv{X}:X\Rightarrow X}
\\[2ex]
\inference{\Psi\ni X:=A}
          {\Psi\vdash\unseal{X}:X\Rightarrow A}
\quad
\inference{\Psi\ni X:=A}
          {\Psi\vdash\seal{X}:A\Rightarrow X}
\\[2ex]
\inference{\Psi\vdash c:A'\Rightarrow A \qquad
           \Psi\vdash d:B\Rightarrow B'}
          {\Psi\vdash c\func d:(A\func B)\Rightarrow(A'\func B')}
\\[2ex]
\inference{\Psi,X\vdash c:A\Rightarrow B}
          {\Psi\vdash\forall X.\,c:(\forall X.A)\Rightarrow(\forall X.B)}
\end{gather*}

}
%===============================================================================
\frame{
  \frametitle{Well-formed types and context morphisms}

\begin{columns}[T]
\begin{column}{0.45\textwidth}
\centering
\fbox{$\Psi\vdash A$}

{\scriptsize
\begin{gather*}
\inference{\Psi\ni X}{\Psi\vdash X}
\quad
\inference{}{\Psi\vdash\iota}
\\[2ex]
\inference{\Psi\vdash A\qquad\Psi\vdash B}
          {\Psi\vdash A\func B}
\\[2ex]
\inference{\Psi,X\vdash A}
          {\Psi\vdash\forall X.A}
\end{gather*}
}
\end{column}
\begin{column}{0.55\textwidth}
\centering
\fbox{$\Psi\vdash\morph{\Theta}{\delta}$}

{\scriptsize
\[
\inference{\delta^{+}(\Psi)\vdash^{\mathrm r}\Theta
           \qquad \Psi\vdash^{\mathrm s}\delta}
          {\Psi\vdash\morph{\Theta}{\delta}}
\]
\fbox{$\Xi\vdash^{\mathrm r} \Theta$}
\begin{gather*}
\inference{}{\Xi\vdash^{\mathrm r}\emptyset}
\quad
\inference{\Xi\vdash A\qquad\Xi\vdash^{\mathrm r}\Theta}
          {\Xi\vdash^{\mathrm r}\bind{X}{A},\Theta}
\end{gather*}
\fbox{$\Psi\vdash^{\mathrm s} \delta$}
\begin{gather*}
\inference{}{\Psi\vdash^{\mathrm s}\emptyset}
\quad
\inference{\delta(\Psi)\ni X\qquad\Psi\vdash^{\mathrm s}\delta}
          {\Psi\vdash^{\mathrm s}\lockX{X},\delta}
\\[2ex]
\inference{\delta(\Psi)\ni_{\mathrm{lk}} X\qquad\Psi\vdash^{\mathrm s}\delta}
          {\Psi\vdash^{\mathrm s}\unlockX{X},\delta}
\end{gather*}
}
\end{column}
\end{columns}
}
%===============================================================================
\frame{
  \frametitle{Term typing}

\hfill\fbox{$\Psi \mid \Gamma \vdash M : A$}

{\scriptsize
\begin{gather*}
\inference{x:A\in\Gamma}
          {\Psi\mid\Gamma\vdash x:A}
\quad
\inference{}
          {\Psi\mid\Gamma\vdash k:\iota}
\\[2ex]
\inference{\Psi\vdash A \qquad
           \Psi\mid\Gamma,x:A\vdash N:B}
          {\Psi\mid\Gamma\vdash \lam{x\of A}N:A\to B}
\quad
\inference{\Psi\mid\Gamma\vdash L:A\to B \qquad
           \Psi\mid\Gamma\vdash M:A}
          {\Psi\mid\Gamma\vdash \app{L}{M}:B}
\\[2ex]
\inference{\Psi,X\mid\Gamma\vdash N:B}
          {\Psi\mid\Gamma\vdash \Lambda X.\,N:\forall X.B}
\quad
\inference{\Psi\mid\Gamma\vdash L:\forall X.B \qquad \Psi\vdash A}
          {\Psi\mid\Gamma\vdash \tyapp{L}{A}:B[X\mapsto A]}
\end{gather*}
}
\[
\frac{
  \begin{gathered}
  \Psi \vdash \Delta
  \qquad
  \intc{\Delta}{\Psi}\mid\emptyset\vdash M:A
  \\[1ex]
  \fce{\Delta}{\Psi}\vdash c:A\Rightarrow B
  \qquad
  \Psi\vdash B
  \end{gathered}
}{
  \Psi\mid\Gamma\vdash\boundary{M}{\Delta}{c}:B
}
\]

}
%===============================================================================
\frame{
  \frametitle{Values, Inert and Active Conversions}

{\small
\[
\begin{array}{llcl}
\text{values} & V,W &::=& k \mid \lam{x\of A}N \mid \Lambda X.\,V
       \mid \boundary{V}{\Delta}{c^i} \\[1ex]
\text{inert}& c^i &::=& \idconv{X} \mid \seal{X}
              \mid c\func d \mid \forall X.\,c \\[1ex]
\text{active}& c^a &::=& \idconv{\iota} \mid \unseal{X}
\end{array}
\]

\vspace{18pt}

\begin{center}
A boundary is a value exactly when its interior is a value\\
and its conversion face has inert form $c^i$.
\end{center}
}
}
%===============================================================================
\frame{
\frametitle{Reduction: elimination rules}

\hfill\fbox{$\Psi \vdash M \reduce N$}

\begin{align*}
\Psi \vdash (\lam{x\of A}N)\,W
  &\reduce N[x\by W]
\\[2ex]
\Psi \vdash \boundary{V}{\Delta}{c\func d}\,W
  &\reduce \boundary{V\,\boundary{W}{\dual{\Delta}}{c}}
                    {\Delta}{d}
\\[2ex]
\Psi \vdash \tyapp{(\Lam{X}V_B)}{A}
  &\reduce \boundary{V}{\morph{\bind{X}{A}}{\emptyset}}{\unsealat{X}{B}}
\\[2ex]
\Psi \vdash \tyapp{\boundary{V}{\morph{\Theta}{\delta}}{\forall X.\,c}}{A}
&\reduce \boundary{\tyapp{V}{X}}{\morph{\bind{X}{A},\Theta}{\delta}}
                  {\unsealatc{X}{c}}
\end{align*}

}
%===============================================================================
\frame{
\frametitle{Reduction: active conversions and congruence}

\hfill\fbox{$\Psi \vdash M \reduce N$}

\begin{align*}
\Psi \vdash \boundary{k}{\Delta}{\idconv{\iota}}
  &\reduce k
\\[2ex]
\Psi \vdash
\boundary{\boundary{V}{\Delta_1}{\seal{X}}}
          {\Delta_2}{\unseal{X}}
&\reduce
\boundary{\boundary{V}{\Delta_1\scopemove\Delta_2}{\idc{A}}}
          {\rewind{\Delta_2}}{\idc{A}}
&& \text{if } \Psi \ni X:=A
\\[2ex]
\Psi \vdash
\boundary{\boundary{V}{\Delta_1}{\idconv{X}}}
          {\Delta_2}{\unseal{X}}
&\reduce
\boundary{\boundary{V}{\Delta_1\scopemove\Delta_2}{\unseal{X}}}
          {\rewind{\Delta_2}}{\idc{A}}
&&\text{if }\Psi\ni X:=A
\end{align*}

{\footnotesize
\begin{align*}
\Psi \vdash L\,M
  &\reduce L'\,M
  &&\text{if }\Psi \vdash L\reduce L'
  \\
\Psi \vdash V\,M
  &\reduce V\,M'
  &&\text{if }\Psi \vdash M\reduce M'
  \\
\Psi \vdash \tyapp{L}{A}
  &\reduce \tyapp{L'}{A}
  &&\text{if }\Psi \vdash L\reduce L'
  \\
\Psi \vdash \Lam{X}N
  &\reduce \Lam{X}N'
  &&\text{if }\Psi,X\vdash N\reduce N'
  \\
\Psi \vdash \boundary{M}{\Delta}{c}
  &\reduce \boundary{M'}{\Delta}{c}
  &&\text{if }\intc{\Delta}{\Psi}\vdash M\reduce M'
\end{align*}
}
}
%===============================================================================
\frame{
\frametitle{Progress}

\[
\begin{aligned}
&\text{If }\Psi\mid\emptyset\vdash M:A \\
&\text{then } M \text{ is a value or } \exists M'.\ \Psi\vdash M\reduce M'.
\end{aligned}
\]

\vspace{24pt}

\begin{center}
Every closed, well-typed term is a value or can take a reduction step.
\end{center}
}
%===============================================================================
\frame{
\frametitle{Preservation}

\[
\begin{aligned}
&\text{If }\Psi\mid\emptyset\vdash M:A
  \text{ and }\Psi\vdash M\reduce M' \\
&\text{then } \Psi\mid\emptyset\vdash M':A.
\end{aligned}
\]

\vspace{24pt}

\begin{center}
A reduction step preserves the type of a closed term.
\end{center}
}
%===============================================================================
\frame{
\frametitle{Determinism}

\[
\begin{aligned}
&\text{If }\Psi\vdash M\reduce M_1
  \text{ and }\Psi\vdash M\reduce M_2 \\
&\text{then } M_1=M_2.
\end{aligned}
\]

\vspace{24pt}

\begin{center}
Every term has at most one immediate reduct.
\end{center}
}

\end{document}
